% ════════════════════════════════════════════════════════════════
%  CONCLUSION GÉNÉRALE ET PERSPECTIVES
% ════════════════════════════════════════════════════════════════

\chapter*{Conclusion générale et perspectives}
\addcontentsline{toc}{chapter}{Conclusion générale et perspectives}

\section*{Bilan du projet}
\addcontentsline{toc}{section}{Bilan du projet}

Ce projet a permis de concevoir et de réaliser une plateforme numérique de
recrutement adaptée à \textbf{Schulte Automotive Tunisia}. Le système
répond au besoin principal de l'entreprise : remplacer un processus
majoritairement manuel par un processus structuré, traçable et plus
transparent.

La solution livrée couvre l'ensemble du cycle : publication des offres,
candidature, suivi RH, planification des entretiens et notifications.
Elle apporte un bénéfice direct pour les trois profils d'utilisateurs :
le candidat, le responsable RH et l'administrateur.

Sur le plan technique, la valeur produite dépasse une application CRUD :
la plateforme orchestre des règles métier multi-acteurs, des contrôles
d'accès par rôle et par site, ainsi qu'une synchronisation entre API,
données persistées et événements temps réel.

Au niveau méthodologique, l'organisation en releases Scrum a permis une
progression maîtrisée, avec des livraisons régulières et des ajustements
guidés par les retours terrain.

\section*{Limites}
\addcontentsline{toc}{section}{Limites}

Les limites identifiées relèvent principalement du périmètre PFE :
les validations ont été réalisées dans un contexte local et contrôlé, et
le déploiement à grande échelle n'a pas été l'objectif de cette version.
L'analyse assistée des candidatures reste un outil d'appui à la décision,
la décision finale restant sous contrôle humain pour garantir l'équité.

\section*{Améliorations futures}
\addcontentsline{toc}{section}{Améliorations futures}

Plusieurs évolutions peuvent être envisagées pour prolonger ce travail :
\begin{itemize}
    \item renforcer l'industrialisation (CI/CD, sauvegardes et monitoring) ;
    \item enrichir l'application mobile avec des notifications push natives ;
    \item étendre les intégrations (emails transactionnels avancés,
    synchronisation calendrier plus complète) ;
    \item préparer un déploiement multi-filiales avec paramétrage par pays
    et par langue.
\end{itemize}

En synthèse, ce projet constitue une base solide et exploitable à court
terme. Les perspectives identifiées montrent que la plateforme peut évoluer
vers une solution RH plus complète, capable d'accompagner durablement la
transformation numérique du recrutement.
